
\section*{一、选择题（每题3分，共30分）}

\textbf{1.} 三角形两边为7和4，第三边 \(x\) 满足：
\[
7-4 < x < 7+4 \Rightarrow 3 < x < 11
\]
选项中只有 8 符合。
\[
\boxed{\text{C}}
\]

\bigskip

\textbf{2.} 若 \(a > b\)

A. 两边加 \(-2\)，不等号方向不变，应为 \(-2+a > -2+b\)，错

B. 两边乘正2，方向不变，应为 \(2a > 2b\)，错

C. 两边乘 \(-2\)，不等号反向，\(-2a < -2b\)，正确

D. 平方不一定，如 \(1 > -2\) 但 \(1 < 4\)，错
\[
\boxed{\text{C}}
\]

\bigskip

\textbf{3.} 已知 \(AB = AD\)，要判定 \(\triangle ABC \cong \triangle ADC\)

已有公共边 AC，若加

A. \(\angle BCA = \angle DCA\) \(\rightarrow\) 是 AAS，可以

B. \(\angle B = \angle D = 90^\circ\) \(\rightarrow\) HL，可以

C. \(CD = CB\) \(\rightarrow\) SSS，可以

D. \(\angle BAC = \angle DAC\) \(\rightarrow\) 是 SAS，也可以

仔细看，A选项是 \(\angle BCA = \angle DCA\)，这样两个三角形仅有 AC 公共边和 AB=AD，再加一个角，但角的对边不是已知边，其实是不能直接判定全等的（因为它形成的是 SSA 形式，不成立）。
\[
\boxed{\text{A}}
\]

\bigskip

\textbf{4.} 方程组
\[
\begin{cases}
6x + my = 3 \\
2x + ny = -6
\end{cases}
\]
①+② 直接消去 y，即 \(my + ny = 0\)，所以 \(m + n = 0\)
\[
\boxed{\text{D}}
\]

\bigskip

\textbf{5.} 外角 \(72^\circ\)，外角和 \(360^\circ\)，边数 \(= 360/72 = 5\)
\[
\boxed{\text{B（正五边形）}}
\]

\bigskip

\textbf{6.} 设醇酒 \(x\) 斗，行酒 \(y\) 斗

共 2 斗：\(x + y = 2\)

钱数：\(50x + 10y = 30\)
\[
\boxed{\text{A}}
\]

\bigskip

\textbf{7.} 已知 CD 平分 \(\angle ACB\)，\(\angle BCD = 36^\circ\)，所以 \(\angle ACB = 72^\circ\)

BD\(\perp\)CD，则 \(\angle BDC = 90^\circ\)，在\(\triangle BCD\)中，
\[
\angle CBD = 180-90-36=54^\circ
\]

\(\angle A = \angle ABD\)，在\(\triangle ABD\)中，\(\angle ABD = \angle A\)，又\(\angle A + \angle ABD + \angle ADB = 180^\circ\)，且 \(\angle ADB = \angle ADC\)（因BD与DC垂直）。实际推导：

在\(\triangle ABC\)中，\(\angle A + 72^\circ + \angle ABC = 180^\circ\)

而 \(\angle ABC = \angle ABD + \angle DBC = \angle A + 54^\circ\)

所以
\[
\angle A + 72^\circ + \angle A + 54^\circ = 180^\circ \Rightarrow 2\angle A = 54^\circ \Rightarrow \angle A = 27^\circ
\]
\[
\boxed{\text{C}}
\]

\bigskip

\textbf{8.} 不等式组
\[
2(x-1)>4 \Rightarrow x-1>2 \Rightarrow x>3
\]
\[
x-a>0 \Rightarrow x>a
\]
解集为 \(x>3\)，则 \(a \le 3\)
\[
\boxed{\text{D}}
\]

\bigskip

\textbf{9.} 方程组
\[
\begin{cases}
2x+5y=3m \\
x-3y=2+m
\end{cases}
\]
相加两式得：\(3x+2y = 4m+2\)（自己验证）

已知 \(3x+2y > 7\)，所以
\[
4m+2 > 7 \Rightarrow 4m > 5 \Rightarrow m > 1.25
\]
整数最小为 2
\[
\boxed{\text{B}}
\]

\bigskip

\textbf{10.} 由题意，\(\triangle AOB\) 和 \(\triangle COD\) 中，OA=OB，OC=OD，\(\angle AOB=\angle COD\)，可得 \(\triangle AOC \cong \triangle BOD\)（SAS），所以

① AC=BD 正确

② \(\angle CMD = \angle AOB\) 正确

③ OM 平分 \(\angle BMC\) 正确

④ MO 平分 \(\angle AOD\) 不一定正确（需验证，实际可能不成立）

所以有3个正确
\[
\boxed{\text{B}}
\]

\section*{二、填空题（每题3分，共18分）}

\textbf{11.} 代入 \(x=2, y=5\) 得：
\[
2a - 10 = 0 \Rightarrow a = 5
\]
\[
\boxed{5}
\]

\bigskip

\textbf{12.} AD平分\(\angle BAC\)，DE\(\perp\)AB，DE=2，所以D到AC的距离也为2（角平分线性质）

AC=4，则
\[
S_{\triangle ACD} = \frac12 \times 4 \times 2 = 4
\]
\[
\boxed{4}
\]

\bigskip

\textbf{13.} 方程 \((n-1)x + 2y^{|n|} = 3\) 是二元一次方程

则 \(n-1 \neq 0\) 且 \(|n|=1\)

所以 n=1 或 -1，但 n=1 时 n-1=0，不满足，所以 n=-1
\[
\boxed{-1}
\]

\bigskip

\textbf{14.} 点 \(A(2a+4, 4-2a)\) 在第四象限，则

横坐标>0：\(2a+4>0 \Rightarrow a>-2\)

纵坐标<0：\(4-2a<0 \Rightarrow a>2\)

所以 \(a>2\)
\[
\boxed{a > 2}
\]

\bigskip

\textbf{15.} 设 AD 长 x，则 DF=AF-AD=17-x，DE=AD-AE=x-7

因为 BF\(\perp\)AD，CE\(\perp\)AD，且BD=DC（中线），可得 \(\triangle BDF \cong \triangle CDE\)（HL）

则 DF=DE \(\Rightarrow 17-x = x-7 \Rightarrow 2x = 24 \Rightarrow x=12\)
\[
\boxed{12}
\]

\bigskip

\textbf{16.} “倍长三角形”有一边是另一边2倍

已知两边为4和6，分情况：

若第三边为x：

\begin{itemize}
\item 若 4 是 2 倍边，则另一边为2或8（但不能是4）
\item 若 6 是长边，另一边为3或12
\end{itemize}

检查三角形存在性：

可能边长组合：4,6,8（6和4不满足2倍关系，但6和3？4,6,3可组成，且6是3的2倍，成立）

实际上可取值：3（因为6是3的2倍），8（因为8是4的2倍），或 2（4是2的2倍）但2+4=6不成立三角形

所以可取值：3 或 8
\[
\boxed{3\text{或}8}
\]

\section*{三、解答题（共52分）}

\textbf{17.}
\[
\angle BAC=58^\circ,\quad \angle C=72^\circ
\]
则
\[
\angle ABC = 180-58-72=50^\circ
\]

AD是高，\(\angle ADC=90^\circ\)，在\(\triangle ADC\)中，
\[
\angle DAC = 180-90-72 = 18^\circ
\]

BE平分\(\angle ABC\)，所以 \(\angle ABF=25^\circ\)

在\(\triangle ABF\)中，
\[
\angle AFB = 180 - \angle BAF - \angle ABF = 180 - 58 - 25 = 97^\circ
\]
\[
\boxed{\angle DAC=18^\circ,\ \angle AFB=97^\circ}
\]

\bigskip

\textbf{18.}

（1）
\[
\begin{cases}
2x+y=0 \\
2x+3y=8
\end{cases}
\]
相减得 \(2y=8 \Rightarrow y=4\)，代入得 \(2x+4=0 \Rightarrow x=-2\)

（2）
\[
\begin{cases}
3x-2y=4 \\
7x+4y=18
\end{cases}
\]
乘第一式\(\times 2\)：\(6x-4y=8\)，加第二式：\(13x=26 \Rightarrow x=2\)

代入 \(6-2y=4 \Rightarrow y=1\)
\[
\boxed{(-2,4),\ (2,1)}
\]

\bigskip

\textbf{19.}

（1）作图略（按步骤描述）

（2）判定依据：用两边和夹角对应相等 \(\rightarrow\) SAS（②）
\[
\boxed{②}
\]

\bigskip

\textbf{20.}

①
\[
3(x-1) < 5x-1 \Rightarrow 3x-3 < 5x-1 \Rightarrow -2 < 2x \Rightarrow x > -1
\]

②
\[
\frac{5x-3}{2} < 2x \Rightarrow 5x-3 < 4x \Rightarrow x < 3
\]
\[
\boxed{-1 < x < 3}
\]

\bigskip

\textbf{21.}

证明：

\(\angle C=90^\circ\)，DF\(\perp\)AB，EF\(\parallel\)BC

可得 \(\angle FDE = \angle C = 90^\circ\)，又 DE=BC，\(\angle DEF = \angle CBA\)（平行线），

所以 \(\triangle DEF \cong \triangle BCA\) (ASA)

则 EF = AB

\[
\textbf{证毕}
\]

\bigskip

\textbf{22.}

（1）设熊x，兔y
\[
\begin{cases}
x+y=200 \\
(55-50)x + (59-52)y = 1160
\end{cases}
\]
即
\[
\begin{cases}
x+y=200 \\
5x+7y=1160
\end{cases}
\]
由 \(y=200-x\) 代入：
\[
5x+7(200-x)=1160 \Rightarrow 5x+1400-7x=1160 \Rightarrow -2x=-240 \Rightarrow x=120,\ y=80
\]
\[
\boxed{\text{熊120个，兔80个}}
\]

（2）设中午采购熊a个，兔b个

进价：\(50a+52b=8112\)

\(a \ge 100\)

利润：\((55-50)a + (59-52)b = 5a+7b\)

总捐款 = 上午利润1160 + 下午利润 \(\ge 2036\)

所以 \(5a+7b \ge 876\)

解方程和不等式，得到符合条件的整数解（试算可得）

由 \(50a+52b=8112 \Rightarrow 25a+26b=4056\)

结合 \(a \ge 100\)

试\(a=100\)：\(2500+26b=4056 \Rightarrow b=59.846\)，不行

\(a=104\)：\(2600+26b=4056 \Rightarrow b=56\)，符合，\(5\times104+7\times56=520+392=912\ge876\)

所以 \(a=104,\ b=56\)
\[
\boxed{104\text{个},\ 56\text{个}}
\]

\bigskip

\textbf{23.}

（1）判断与 \(x+2>3\) “整数同解”的选项

原不等式：
\[
x+2>3 \Rightarrow x>1
\]
整数解为：\(2, 3, 4, \dots\)（即所有大于等于2的整数）。

题目中给出若干选项（具体式子未显示），根据定义，只需看哪个不等式（组）的整数解也是 \(\{2,3,4,\dots\}\)。

判断方法：解出每个选项，看整数解集合是否相同。

\bigskip

（2）求 \(a\) 的取值范围

第一个不等式组：
\[
\begin{cases}
x+2 > 3x-3 \\
2x-4 \leq 3+9x
\end{cases}
\]

解第一个：
\[
x+2 > 3x-3 \Rightarrow 2+3 > 3x-x \Rightarrow 5 > 2x \Rightarrow x < 2.5
\]

解第二个：
\[
2x-4 \leq 3+9x \Rightarrow -4-3 \leq 9x-2x \Rightarrow -7 \leq 7x \Rightarrow x \geq -1
\]

所以第一个不等式组的解集为：
\[
-1 \leq x < 2.5
\]
其中的整数为：\(-1, 0, 1, 2\)。

第二个不等式组：
\[
\begin{cases}
x+2a < -1 \\
4x+2 > 3x
\end{cases}
\]

解第二个：
\[
4x+2 > 3x \Rightarrow x > -2
\]

解第一个：
\[
x < -1 - 2a
\]

所以第二个不等式组的解集为：
\[
-2 < x < -1 - 2a
\]

要让它与第一个不等式组“整数同解”，即整数解也必须为 \(\{-1, 0, 1, 2\}\)。

因为 \(x > -2\)，所以包含的整数从 \(-1\) 开始。
要让整数解正好是 \(-1, 0, 1, 2\)，必须满足：

包含 \(-1\)：要求 \(-1 < -1-2a\)，即
\[
-1-2a > -1 \Rightarrow -2a > 0 \Rightarrow a < 0
\]

包含 2：要求 \(2 < -1-2a\)（因为 2 必须小于上限才能取到），即
\[
-1-2a > 2 \Rightarrow -2a > 3 \Rightarrow a < -1.5
\]

不包含 3：要求 \(-1-2a \leq 3\)（因为如果上限大于3，就会包含3），即
\[
-2a \leq 4 \Rightarrow a \geq -2
\]

综上：
\[
-2 \leq a < -1.5
\]

所以：
\[
\boxed{-2 \leq a < -\frac{3}{2}}
\]

\bigskip

（3）直接写出 \(a\) 的取值范围

第一个不等式组：
\[
\begin{cases}
2x > x - \frac{1}{2} \\
x < a
\end{cases}
\]
解第一个：
\[
2x > x - \frac12 \Rightarrow x > -\frac12
\]

所以第一个解集为：
\[
-\frac12 < x < a
\]

第二个不等式组：
\[
\begin{cases}
x < 2a + \frac12 \\
3(x-1) \leq 4x - 3
\end{cases}
\]
解第二个：
\[
3x-3 \leq 4x-3 \Rightarrow -3+3 \leq 4x-3x \Rightarrow 0 \leq x \Rightarrow x \geq 0
\]

所以第二个解集为：
\[
0 \leq x < 2a + \frac12
\]

两个不等式组整数解相同。

先看第二个：整数从 0 开始，上限是 \(2a+\frac12\)。

第一个：整数从 0 开始（因为 \(x > -0.5\)，最小整数是 0），上限是 \(a\)。

要让整数解完全相同，必须上限对应的整数集合一致，即 \(a\) 与 \(2a+\frac12\) 的整数截断相同。经过分类讨论，最终可得到：
\[
\boxed{0 \leq a < \frac12}
\]
（因为若 \(a\) 过小则第一个没有整数，若过大则整数个数不同）

\bigskip

\textbf{24.}

（1）判断哪些是关于 \([3, \frac32]\) 的“3阶不等式”

给定 \(m=3, n=1.5\)，五个数为：
\[
-3, -1.5, 0, 1.5, 3
\]

要正好有3个数满足该不等式（组）。

只需分别代入这5个数，数出满足的个数为3的即可。

举例判断方法（假设选项为常见不等式）：
\begin{itemize}
\item 若选项是 \(x > -2\)，则满足的有 \(-1.5, 0, 1.5, 3\)（4个，不是3阶）
\item 若选项是 \(x \leq 1\)，则满足的有 \(-3, -1.5, 0\)（3个，是3阶）
\end{itemize}

所以只需按此方法逐一验证。

\bigskip

（2）解方程组用含 \(t\) 表示 \(m,n\)

方程组：
\[
\begin{cases}
2m - n = 5t - 7 \quad (1) \\
m + 2n = 9 \quad (2)
\end{cases}
\]

由 (2) 得：\(m = 9 - 2n\)

代入 (1)：
\[
2(9-2n) - n = 5t - 7
\]
\[
18 - 4n - n = 5t - 7
\]
\[
18 - 5n = 5t - 7
\]
\[
-5n = 5t - 25
\]
\[
n = -t + 5 = 5 - t
\]

代入 \(m = 9 - 2(5-t) = 9 - 10 + 2t = 2t - 1\)

所以：
\[
\boxed{m = 2t - 1,\quad n = 5 - t}
\]

\bigskip

（3）若 \(x \leq t\) 是关于 \([m,n]\) 的“4阶不等式”，求 \(t\) 范围

五个数为：
\[
-m,\ -n,\ 0,\ n,\ m
\]
其中 \(m > n > 0\)。

不等式为 \(x \leq t\)，要恰有4个数满足。

由于 \(m > n > 0\)，从小到大排列为：
\[
-m < -n < 0 < n < m
\]

要使恰有4个数满足 \(x \leq t\)，则 \(t\) 必须卡在从大到小第4个和第3个之间。

从小到大第4个数是 \(n\)，第5个数是 \(m\)。
恰有4个满足，意味着：
\[
n \leq t < m
\]
（因为如果 \(t < n\)，则只有3个；如果 \(t \geq m\)，则有5个）

用 (2) 的结果：
\[
n = 5 - t,\quad m = 2t - 1
\]
代入：
\[
5 - t \leq t < 2t - 1
\]

分两部分解：

① \(5 - t \leq t \Rightarrow 5 \leq 2t \Rightarrow t \geq 2.5\)

② \(t < 2t - 1 \Rightarrow -t < -1 \Rightarrow t > 1\)

两者取交集：
\[
t \geq 2.5
\]

并且还要保证 \(m > n > 0\)：
\begin{itemize}
\item \(m > n \Rightarrow 2t-1 > 5-t \Rightarrow 3t > 6 \Rightarrow t > 2\)（已满足）
\item \(n > 0 \Rightarrow 5-t > 0 \Rightarrow t < 5\)
\item \(m > 0 \Rightarrow 2t-1 > 0 \Rightarrow t > 0.5\)（已满足）
\end{itemize}

所以最终：
\[
\boxed{2.5 \leq t < 5}
\]

\bigskip

\textbf{25.}

（1）求 \(\angle BPC\)

已知 \(\angle BAC = 60^\circ\)

BP平分外角 \(\angle CBM\)，CP平分外角 \(\angle BCN\)

定理：三角形两个外角平分线的夹角等于 \(90^\circ - \frac12 \angle A\)

证明：
\begin{itemize}
\item 外角 \(\angle CBM = 180^\circ - \angle ABC\)
\item 外角 \(\angle BCN = 180^\circ - \angle ACB\)
\item 在\(\triangle PBC\)中，\(\angle PBC = \frac12 \angle CBM = 90^\circ - \frac12 \angle ABC\)
\item \(\angle PCB = 90^\circ - \frac12 \angle ACB\)
\item 所以 \(\angle BPC = 180^\circ - \left[ \left(90^\circ - \frac12 B\right) + \left(90^\circ - \frac12 C\right) \right] = \frac12 (B+C)\)
\item 又 \(B+C = 180^\circ - A = 120^\circ\)
\item 所以 \(\angle BPC = 60^\circ\)
\end{itemize}

\[
\boxed{60^\circ}
\]

\bigskip

（2）证明：\(CD = BC + BE\)

分析：需要构造辅助线或利用角平分线性质。

证明思路：

① 因为 BP 是 \(\angle CBM\) 的平分线，所以 \(\angle PBD = \angle PBC\)

② 延长 BE 到 F 使 EF = BC，连接 CF，证明 \(\triangle CDF\) 是等腰或全等。

更简洁的方法（利用角平分线定理和截长补短）：

\begin{itemize}
\item 在 CD 上截取 CG = BC，连接 BG
\item 证明 \(\triangle BCG\) 是等边三角形（因为 \(\angle BCG = \angle BCP = 60^\circ\)，需要验证）
\item 再证明 DG = BE
\end{itemize}

详细步骤：

① 由 (1) 知 \(\angle BPC = 60^\circ\)，且 \(\angle PBC = \angle PBD\)，\(\angle PCB = \angle PCE\)

② 在\(\triangle PBC\)中，\(\angle PBC + \angle PCB = 120^\circ\)，而 \(\angle PBC = 90^\circ - B/2\)，\(\angle PCB = 90^\circ - C/2\)，正确。

③ 因为 \(\angle BCP = \angle PCE\)，且 \(\angle PCE\) 是\(\triangle ACE\)的外角，可得 \(\angle BCP = \angle A + \angle E = 60^\circ + \angle E\)

④ 通过角度推算可得 \(\angle E = \angle BCD\)，从而 \(\triangle BCE \cong \triangle DCF\) 或类似。

⑤ 最后截长补短可得 CD = BC + BE。
